\documentclass[11pt]{article}
\usepackage[margin=0.9in]{geometry}
\usepackage{amsmath,amssymb,amsthm,booktabs,array}
\usepackage[T1]{fontenc}
\usepackage{hyperref,xurl}
\hypersetup{colorlinks=true,urlcolor=blue,linkcolor=blue}
\newtheorem{proposition}{Proposition}
\newcommand{\HomR}{\operatorname{Hom}_R}
\newcommand{\Span}{\operatorname{span}}
\newcommand{\src}[1]{\nolinkurl{#1}}
\title{Lean research notes: reusable mathematical reductions}
\author{Principia contributors}
\date{21 November 2026}
\begin{document}
\maketitle
\begin{abstract}
We record checked replacements in Alp\"oge's formal development: separation
of an isotypic module by equivariant projections, extension of subspace
isomorphisms, projective splitting, torsion finiteness, and a common
polynomial certificate for three projective addition identities.
The replacements preserve the existing public theorem types.
The mathematical families save 867 physical source lines, including their
new shared helper. A further stabilizer abstraction saves six proof-body
lines. Expanded-statement cleanup is counted separately.
\end{abstract}

\section{An isotypic module is detected by its simple type}
Let $R$ be a ring, $S$ a simple left $R$-module, and $M$ another module.
Suppose
\[
 M=\sum_{f\in\HomR(S,M)}\operatorname{im} f.
\]
\begin{proposition}
The module $M$ is semisimple and isotypic of type $S$. For every
$0\ne v\in M$, there is $p\in\HomR(M,S)$ with $p(v)\ne0$.
Every nonzero submodule of $M$ contains an embedded copy of $S$.
\end{proposition}
\begin{proof}
A nonzero map from $S$ is injective, so each nonzero image is a copy of $S$.
The spanning condition therefore makes the isotypic component of type $S$
equal to $M$, and expresses $M$ as a sum of semisimple submodules.
Mathlib's isotypic decomposition gives
$M\simeq\bigoplus_{i\in I} S$. A nonzero vector has a nonzero coordinate;
the corresponding coordinate map is $p$. Finally, a nonzero submodule of a
semisimple module contains a simple submodule, and isotypy identifies that
submodule with $S$.
\end{proof}

For representations over a field $k$, apply this to the monoid algebra
$R=k[G]$. The resulting map $p$ is an intertwiner. If a subgroup $U$
fixes $v$, then $p(v)$ is a nonzero $U$-fixed vector in $S$.
The algebraic statement requires no topology, compactness, smoothness or
finite-index hypothesis.

This replaces compact-group averaging in \src{HasFixedVector}, finite-sum
induction in \src{FixedVectorsNeBotOfMemSumIntertwiningRanges}, and repeated
projection constructions in two other consumers. The same lemma reduces
\src{GlobalTheory.IsoSep.main} from a 139-line module to 16 lines: its
additional isomorphism premise is retained in the public interface but is
unnecessary for the new proof. The reusable implementation is
\src{Lemmata/IsotypicGeneration.lean}.

\section{One cubic certificate replaces three expanded certificates}
Write a projective point as $P=(x,y,z)$, and let $F(P)=0$ be its
homogeneous Weierstrass equation. Define
\begin{align*}
r(P)={}&(a_1^2+a_2)x^3+3xy(a_1x+y+a_3z)\\
 &+3xz\big((a_1a_3+a_4)x+(a_3^2+3a_6)z\big)+b_8z^3,
\end{align*}
where $b_8=a_1^2a_6+4a_2a_6-a_1a_3a_4+a_2a_3^2-a_4^2$.
The three cross-coordinate compatibility identities for the two projective
addition formulas each admit certificate multipliers
\[
 -C(P,Q)r(Q),\qquad C(P,Q)r(P),
\]
against the two curve equations. Thus they all factor through the same
alternating residual $r(P)F(Q)-r(Q)F(P)$. Each coefficient $C$ has only a
few bilinear terms. For example, the $X,Z$ identity uses
\[
 C(P,Q)=a_1x_Pz_Q+y_Pz_Q+z_Py_Q+a_3z_Pz_Q.
\]

The original expanded multipliers do not factor directly. The useful step
is to change a certificate $(A,B)$ by
\[
 (A,B)\longmapsto\big(A-HF(Q),\ B+HF(P)\big).
\]
This preserves $AF(P)+BF(Q)$ by commutativity. Choosing this zero relation
before factoring exposes the common cubic. Changing $r$ by a multiple of
$F$ likewise leaves the alternating residual unchanged.

The three original certificate expressions occupy 102, 142 and 223 physical
lines. Their replacements occupy 5, 6 and 7 lines, plus the shared cubic.
The full \src{DiagZEqNegAddY.lean} module falls from 629 to 189 lines,
saving 440 after including the new definition. Lean checks the resulting
polynomial identities over an arbitrary commutative ring; no division or
nonsingularity assumption is added.

\section{Extension, splitting, and duality replace coordinate constructions}
An isomorphism between subspaces of a finite-dimensional vector space
extends to an automorphism of the ambient space. Applying Mathlib's extension
theorem replaces separate complement choices and block-product identifications
in three tuple/subspace arguments.

If $f:V\to k^n$ is surjective, projectivity supplies a linear section $s$.
The decomposition $V=\ker f\oplus\operatorname{im}s$ follows from the
bijectivity of
\[
 \ker f\oplus k^n\longrightarrow V,\qquad (u,y)\longmapsto u+s(y).
\]
This avoids constructing each lift separately and writing two explicit
inverse maps. Existing public basis conclusions and their hypotheses are
preserved.

For the dual counting arguments, independent coordinate functionals make
the dual map injective. Finite-dimensional duality then gives surjectivity
of the original map, replacing kernel-dimension, projection, and annihilator
arguments. Across four modules these changes remove 165 physical lines;
the changed proof bodies fall from 281 to 119 nonempty lines. No new shared
module is needed for this family.

\section{Torsion, modularity, and stabilizers}
The module named \src{QZ} is definitionally $\mathbb Q/\mathbb Z$, represented
by Mathlib's additive circle of period $1$. Consequently, for $m\ne0$,
\[
 \{x\in\mathbb Q/\mathbb Z:mx=0\}
\]
is finite by \src{AddCircle.finite_torsion}. This replaces a 27-line
residue enumeration by one line.

For subgroups $U\trianglelefteq G$ and $U\le X$, Dedekind's identity yields
$X\cap(HU)=(X\cap H)U$. Converting products to joins gives the required
subgroup equality directly. Finite-index inheritance handles the related
nonzero-index claims. This family saves 31 physical lines.

For a representation of a finite product group, the stabilizer of a vector
is already a subgroup supplied by Mathlib. The theorem that a product
subgroup is generated by its coordinate embeddings then gives the
fixed-vector equivalence. This replaces the private helper's manual
stabilizer construction, reducing its proof from 11 to 5 lines.

\section{Counts, attribution, and verification}
\begin{center}
\begin{tabular}{lr}
\toprule
Family & Net physical lines removed\\
\midrule
Shared isotypic argument, including all new helper code & 231\\
Common polynomial certificate & 440\\
Subspace extension, splitting, and duality & 165\\
Torsion and subgroup identities & 31\\
\midrule
Mathematical families & 867\\
\bottomrule
\end{tabular}
\end{center}
These are source-line measurements, including imports, comments and spacing;
they are not percentages of mathematical difficulty. The separate fixed-vector
helper saves six proof-body lines. Its surrounding expanded public statement
also contracts substantially: that entire module falls from 8,029 to 33 lines.
Most of that latter saving is inference/statement cleanup and is excluded from
the 867-line table.

The concise fixed-vector interface follows
\href{https://github.com/leanprover-community/riemann-hypothesis/pull/44}{mmmbeer's PR 44},
pinned at \src{c9b3e57f0fe987012d32c2523635387bfe56d5a5}. Its general
product-subgroup argument was already present in the original archive;
we do not count that existing abstraction as new work. The proofs build on
Alp\"oge's archive and the cited Mathlib theorems.

All changed modules have focused native Lean checks with Lean 4.33.1 and
Mathlib \src{0df444a360eaa60ab8c11dca51a86af692955474}.
The changed audited endpoints report only propositional extensionality,
classical choice and quotient soundness. Selected public types are compared
directly with their original compiled expressions; other unchanged public
headers are preserved byte for byte. Additional empty-kernel replay checks
cover \src{mem_UC} (38,790 declarations) and \src{GlobalTheory.IsoSep.main}
(12,291 declarations), including their transitive proof dependencies.

The repository contains source snapshots, hashes, commands, logs and an
independent review under \src{verification/research/nov21-compaction/}.
These focused results do not supply a full-project Comparator verdict or
replace the central arithmetic/geometric RH argument. The five final
statements, Challenge, Solution, permitted axioms and dependency pins are
preserved.
\end{document}
